\section{One-sided tangent neighborhoods}
\label{sec:tangentneighborhood86}
The preceding curve results admit a finite-pair formulation.  The
base measurement and its direction are fixed, but the second-order
perturbation may vary arbitrarily within a prescribed bound.  This
also permits first-order opening of a missing support.  Write
$\|R\|_{\Sigma}=\sum_j\|R_j\|_{\mathrm{op}}$ for a Hermitian tuple.
All trace and diamond norms in this section are unhalved.

\begin{lemma}[A reference-uniform channel bound]\label{lem:channelnorm86}
For a Hermitian tuple $R$, let
$\mathcal D_R(X)=\sum_j\tr(R_jX)|j\rangle\langle j|$.
Then
\begin{equation}\label{eq:channelnorm86}
 \|\mathcal D_R\|_\diamond\le\|R\|_{\Sigma}.
\end{equation}
The map is Hermiticity preserving, and is trace annihilating when
$\sum_jR_j=0$.  In particular this applies to measurement-channel
differences, derivatives, and Taylor remainders.
\end{lemma}
\begin{proof}
At every reference dimension the output is block diagonal.  In trace
norm duality it therefore suffices to use a block-diagonal test
$Y=\bigoplus_jY_j$, with $\|Y_j\|_{\mathrm{op}}\le1$.  Its adjoint
image is $\sum_jR_j\otimes Y_j$, of norm at most
$\sum_j\|R_j\|_{\mathrm{op}}$.  Duality bounds the amplified map on
\emph{every} trace-class input $X$, not only on states.  Taking the
supremum over reference dimensions proves \eqref{eq:channelnorm86}.
The remaining assertions follow by taking adjoints and traces.
\end{proof}

\begin{lemma}[The one-sided cone and its lineality]\label{lem:onesidedcone86}
Fix a measurement $E$, put $P_j=\operatorname{supp}E_j$ and $Q_j=I-P_j$,
and let $H$ be a Hermitian zero-sum tuple.  It is the right derivative
at zero of a one-sided real-analytic measurement curve if and only if
\begin{equation}\label{eq:onesidedcone86}
 J_j:=Q_jH_jQ_j\succeq0\quad\hbox{for every }j.
\end{equation}
The largest vector subspace of this cone is given by $J_j=0$ for all
$j$.  For every Hermitian zero-sum tuple, not necessarily in the cone,
\begin{equation}\label{eq:fullrange86}
 H\in\operatorname{Ran}\mathcal C_E
 \quad\Longleftrightarrow\quad
 J_j=0\ (1\le j\le k),\qquad
 \Gamma_E(H)\in\mathcal W_E.
\end{equation}
Here $\Gamma_E(H)=\frac i2\sum_j[P_j,H_j]$ and
$\mathcal W_E=\sum_jP_j\mathcal H_dP_j$ are real Hermitian objects.
\end{lemma}
\begin{proof}
Right differentiation of $\langle v,E_j(t)v\rangle\ge0$ for
$v\in\operatorname{Ran}Q_j$ proves necessity.  For sufficiency, for
each nonzero support let $\lambda_j>0$ be the smallest eigenvalue
of $E_j$ on that support.  Choose
\[
 c\ge1+\max_{j:P_j\ne0}2\|H_j\|_{\mathrm{op}}^2/\lambda_j.
\]
For all sufficiently small $s\ge0$, the $P_j$ block of
$E_j+sH_j+cs^2I$ is at least $\lambda_jP_j/2$.
Its Schur complement on $Q_j$ is bounded below by
\[
 sJ_j+s^2\left(c-2\|H_j\|_{\mathrm{op}}^2/\lambda_j\right)Q_j
 \succeq0.
\]
If $P_j=0$, condition \eqref{eq:onesidedcone86} is $H_j\succeq0$
and positivity is immediate.  Thus the rational matrix curve
\begin{equation}\label{eq:conecurve86}
 E_j^{\mathrm c}(s)=\frac{E_j+sH_j+cs^2I}{1+kcs^2}
\end{equation}
is positive, sums exactly to $I$, and has right derivative $H_j$.
This proves sufficiency, including singular $J_j$ with nonzero cross
blocks.  The cone and its negative intersect precisely at $J_j=0$.

Pair $H$ with the complete kernel in
Theorem~\ref{thm:supportkernel84}.  Taking only its independent
missing-support parameters $Z_j$ shows that orthogonality to the
kernel forces $Q_jH_jQ_j=0$; the tuple-mean term vanishes because
$\sum_jH_j=0$.  Once those blocks vanish, the remaining pairing is
$-2\tr(A\Gamma_E(H))$ for every $A\in\mathcal W_E^\perp$.
Self-adjointness of $\mathcal C_E$ proves \eqref{eq:fullrange86}.
\end{proof}

\begin{lemma}[A finite normalized-factor remainder]\label{lem:factorremainder86}
Let $A_j^*A_j=E_j$, $A^*A=I$, and
$A_j^*B_j+B_j^*A_j=H_j$, with $\sum_jH_j=0$.
Set $b=\|B\|_{\mathrm{op}}$, $h=\max_j\|H_j\|_{\mathrm{op}}$ and
\[
 R_s=(I+s^2B^*B)^{-1/2},\qquad
 G_j(s)=R_s(A_j+sB_j)^*(A_j+sB_j)R_s.
\]
This defines an analytic measurement near zero, and, for $0\le s\le1$,
\begin{equation}\label{eq:factorremainder86}
 \|G(s)-E-sH\|_{\Sigma}\le K_b s^2,
 \qquad K_b=kb^2(2+h).
\end{equation}
When $sb\le1$, its stacked factors differ from $A$ in operator norm
by at most $3bs/2$.  If $A^*B=0$, then also
\begin{equation}\label{eq:horizontalpair86}
 D_N(E,G(s))\le\min\{2,2b\sqrt N s\}.
\end{equation}
\end{lemma}
\begin{proof}
The zero tuple sum gives $A^*B+B^*A=0$, so the stacked normalization
is $I+s^2B^*B$.  This matrix is positive definite for all real $s$;
its inverse square root is analytic near zero.  Functional calculus
gives $\|R_s\|\le1$ and $\|R_s-I\|\le s^2b^2/2$.
Expanding the effect as
$R_s(E_j+sH_j+s^2B_j^*B_j)R_s$ bounds its remainder by
$s^2b^2(1+sh)+s^2\|B_j\|^2$.  Sum over $j$, use
$\sum_j\|B_j\|^2\le kb^2$, and obtain \eqref{eq:factorremainder86}.
The factor difference is at most $sb+s^2b^2/2$.
For horizontal $B$, the ordered identities in
Lemma~\ref{lem:horizontalcurve85} give $A(s)^*A'(s)=0$ and
$\|A'(s)\|\le b$.  Copying the classical label into a discarded
environment is an isometric embedding of the stacked factors, so
it preserves these derivative norms.  The orthogonal-slot argument
of that lemma gives \eqref{eq:horizontalpair86}, including bounded
public stopping.  The environment is not supplied to the tester.
\end{proof}

\begin{definition}[A quadratic tangent neighborhood]\label{def:tangenttube86}
Fix $E$, a Hermitian zero-sum $H\ne0$ satisfying
\eqref{eq:onesidedcone86}, and a number $\Lambda\ge0$.
For $s>0$, the neighborhood consists of all legal measurements $F$
with
\begin{equation}\label{eq:tangenttube86}
 \|F-E-sH\|_{\Sigma}\le\Lambda s^2.
\end{equation}
The number $\Lambda$ is a bound on a \emph{finite remainder}, not an
inferred curvature bound.  No curve, or differentiability of $F$ as
a function of $s$, is part of this definition.
\end{definition}

\begin{theorem}[Finite-pair classification in a fixed tangent neighborhood]
\label{thm:tubedichotomy86}
For each fixed $E,H,\Lambda$ in Definition~\ref{def:tangenttube86}
there are $c,C,s_0>0$ such that, for every integer $N\ge1$, every
$0<s\le s_0$, and every legal $F$ satisfying
\eqref{eq:tangenttube86},
\begin{equation}\label{eq:tubedichotomy86}
 c\min\{1,N^\alpha s\}\le D_N(E,F)
       \le C\min\{1,N^\alpha s\},\qquad
 \alpha=\begin{cases}
  1/2,&H\in\operatorname{Ran}\mathcal C_E,\\
  1,&H\notin\operatorname{Ran}\mathcal C_E.
 \end{cases}
\end{equation}
The constants are uniform over the stated remainder neighborhood,
but not over $E,H$ or $\Lambda$.  If some $J_j\ne0$, the linear
lower is attained by repeated product inputs and a classical event.
If all $J_j=0$ but $\Gamma_E(H)\notin\mathcal W_E$, it is attained
by the known-direction correction code on the actual label and a
retained reference, with ideal controls.  Neither construction is
a common unknown-measurement learner.
\end{theorem}

There are two different reasons for a linear scale.  First-order
support opening creates an event impossible at the base measurement.
Within the lineality space of the cone, coherent correction produces
the established Kraus-span metrological mechanism \cite{ZJ84,KL84}.
The result above makes this distinction in support coordinates and
controls finite remainders uniformly at a fixed base and direction.
It is not a uniform arbitrary-pair equivalence for the midpoint
covariance certificate.

\begin{proof}
Put $h=\max_j\|H_j\|_{\mathrm{op}}>0$ and
$h_1=\|H\|_{\Sigma}$.  Choose a unit eigenvector of one $H_j$
with expectation of magnitude $h$.  On that fixed input the event
``label $j$'' has probability gap at least $hs-\Lambda s^2$.
For $s$ sufficiently small this is at least $hs/2$.
Lemma~\ref{lem:bernoulli85} therefore gives a uniform product lower
of order $\min\{1,\sqrt N s\}$, for every $F$ in the neighborhood.
All pairs also satisfy the elementary upper
\begin{equation}\label{eq:tubehybrid86}
 D_N(E,F)\le\min\{2,N(h_1s+\Lambda s^2)\},
\end{equation}
by Lemma~\ref{lem:channelnorm86} and channel telescoping.

Suppose first that $H\in\operatorname{Ran}\mathcal C_E$.
The horizontal solution of Lemma~\ref{lem:horizontalcurve85}
provides factors $B$ with $A^*B=0$ and $T_AB=H$.
Lemma~\ref{lem:factorremainder86} and \eqref{eq:tangenttube86} give
\begin{equation}\label{eq:tubefiniteupper86}
 D_N(E,F)\le\min\{2,2b\sqrt N s+(\Lambda+K_b)Ns^2\}.
\end{equation}
For $x=\sqrt N s\le1$, the right side is at most
$(2b+\Lambda+K_b)x$; for $x\ge1$ it is at most two.
This proves the upper in the first branch without discarding its
finite $Ns^2$ remainder before comparison.

Next suppose some $J_j\ne0$.  It is positive semidefinite; choose a
unit vector $v\in\operatorname{Ran}Q_j$ with
$\lambda=\langle v,J_jv\rangle>0$.
The base probability of label $j$ is zero, while that under $F$ is
at least $\lambda s-\Lambda s^2\ge\lambda s/2$ after decreasing
$s_0$.  Repeating $v$ gives the exact impossible-event bound
\begin{equation}\label{eq:leakagelower86}
 D_N(E,F)\ge2\left[1-(1-\lambda s/2)^N\right]
 \ge2(1-e^{-1})\min\{1,N\lambda s/2\}.
\end{equation}
Here $s_0$ is chosen also so that $\lambda s/2\le1$.
Together with \eqref{eq:tubehybrid86} this proves the linear law.
The vector and the decision event are the same for every allowed $F$.

It remains to consider $J_j=0$ for all $j$ and
$\Gamma:=\Gamma_E(H)\notin\mathcal W_E$.
Use the completed recovery and logical code of
Lemma~\ref{lem:supportcode84}, applied to $\Gamma$.
The normalized first-order factors in Lemma~\ref{lem:curvelogical85}
show algebraically that the corrected derivative $\mathcal D_H$
is $-i[\Gamma_L,\cdot]$.  This calculation needs only $E,H$, not
a smooth factorization of $F$.
For $\Phi_F=\mathcal R\circ(\mathcal M_F\otimes\operatorname{Id})
\circ\mathcal J$, \eqref{eq:tangenttube86} gives
\[
 \|\Phi_F-\operatorname{Id}_2+is[\Gamma_L,\cdot]\|_\diamond
 \le\Lambda s^2.
\]
The integral Taylor remainder for the logical unitary is at most
$2\|\Gamma\|_{\mathrm{op}}^2s^2$.  Consequently
\begin{equation}\label{eq:tubelogical86}
 \|\Phi_F-\operatorname{Ad}(e^{-is\Gamma_L})\|_\diamond
 \le\kappa_{\mathrm{tube}}s^2,
 \qquad \kappa_{\mathrm{tube}}=\Lambda+2\|\Gamma\|_{\mathrm{op}}^2.
\end{equation}
The gap is $\Delta=\|Z\|_{\mathrm{HS}}^2/\tr Z_+>0$, where
$Z=\Gamma-\Pi_{\mathcal W_E}\Gamma$.
For every $m\le N$, equal logical amplitudes and channel telescoping
give a lower bound
$\bigl(2|\sin(ms\Delta/2)|-m\kappa_{\mathrm{tube}}s^2\bigr)_+$.
Choose
\[
 s_0\le\min\{1,(2\Delta)^{-1},
              \Delta/(2\pi\kappa_{\mathrm{tube}})\},\qquad
 m=\min\{N,\lfloor(s\Delta)^{-1}\rfloor\}.
\]
The floor is at least two, and
$\frac12\min\{1,Ns\Delta\}\le x=ms\Delta\le1$.
The sine term is at least $2x/\pi$ and the loss at most $x/(2\pi)$.
In particular the separation is at least
$\min\{1,Ns\Delta\}/(2\pi)$.
The readout may be fixed to the ideal logical pair, so this lower is
uniform over all the permitted $F$ as well.  Its code uses an external
reference of dimension at most $2d$ per call.  Finally
\eqref{eq:tubehybrid86} completes the proof.
\end{proof}

\begin{corollary}[All nonzero one-sided first jets]\label{cor:onesidedcurves86}
Let $E(t)$, $0\le t<a$, be a fixed $C^2$ measurement curve, including its
right derivatives at zero, with $H=E'(0)\ne0$.  Then
\eqref{eq:tubedichotomy86} holds with $s=t$ and curve-dependent
constants on a fixed right interval.  The branch is decided by
\eqref{eq:fullrange86}.  Thus first-order support openings, as well
as tangential second-order rank changes, are included.
\end{corollary}
\begin{proof}
One-sided positivity gives \eqref{eq:onesidedcone86}.
On $[0,a_0]$ set
$L_+=\sup_{0\le u\le a_0}\sum_j\|E_j''(u)\|_{\mathrm{op}}$.
Taylor's formula gives \eqref{eq:tangenttube86} with
$\Lambda=L_+/2$.  Apply Theorem~\ref{thm:tubedichotomy86}.
This does not extend the two-sided tangent lemma by dropping one of
its hypotheses; it uses the larger cone of Lemma~\ref{lem:onesidedcone86}.
\end{proof}

\begin{corollary}[Certified controls on the whole neighborhood]
\label{cor:tubecontrols86}
In the tangential linear branch, Theorem~\ref{thm:controlstability85}
holds uniformly over \eqref{eq:tangenttube86} with $t=s$ and
$\kappa=\kappa_{\mathrm{tube}}$ from \eqref{eq:tubelogical86}.
In particular the same initial, cycle and final error allowances
with $z=\min\{1,Ns\Delta\}$ preserve separation at least $z/(2\pi)$.
\end{corollary}
\begin{proof}
Equation~\eqref{eq:tubelogical86} supplies the only curvature input
used in that theorem's telescoping proof.  Pay it once, and pay the
certified implementation errors under both hypotheses exactly as
in \eqref{eq:controlledger85}.  The ideal readout depends on
$E,H,s,m$, not on the unknown remainder within the neighborhood.
The precision allowances remain angle dependent, not fixed-noise
robustness or a control-synthesis guarantee.
\end{proof}
